\documentclass[11pt,a4paper]{article}
\usepackage[margin=2.5cm]{geometry}
\usepackage{amsmath,amssymb,amsthm}
\usepackage{booktabs}
\usepackage{graphicx}
\usepackage{hyperref}
\usepackage{float}
\usepackage{url}
\usepackage{microtype}
\usepackage{enumitem}

\newcommand{\rad}{\operatorname{rad}}
\newcommand{\Jac}{\operatorname{Jac}}
\newcommand{\Cond}{\operatorname{Cond}}
\newcommand{\End}{\operatorname{End}}
\newcommand{\Res}{\operatorname{Res}}
\newcommand{\Tr}{\operatorname{Tr}}
\newcommand{\Frob}{\operatorname{Frob}}
\newcommand{\Aut}{\operatorname{Aut}}
\newcommand{\GSp}{\mathrm{GSp}}
\newcommand{\GL}{\mathrm{GL}}
\newcommand{\USp}{\mathrm{USp}}
\newcommand{\Gm}{\mathbb{G}_{\mathrm{m}}}
\newcommand{\Ga}{\mathbb{G}_{\mathrm{a}}}
\newcommand{\Qbar}{\overline{\mathbb{Q}}}
\newcommand{\Ql}{\mathbb{Q}_\ell}

\newtheorem{theorem}{Theorem}[section]
\newtheorem{proposition}[theorem]{Proposition}
\newtheorem{corollary}[theorem]{Corollary}
\newtheorem{conjecture}[theorem]{Conjecture}
\newtheorem{lemma}[theorem]{Lemma}
\theoremstyle{definition}
\newtheorem{definition}[theorem]{Definition}
\theoremstyle{remark}
\newtheorem{remark}[theorem]{Remark}

\title{Palindromic Genus-$2$ Curves as Weil Restrictions:\\
  Trace Vanishing, Asai $L$-Functions,\\
  and Modularity via $\mathbb{Q}(i)$}
\author{Chen Ruqing\\[4pt]
  \small GUT Geoservice Inc., Montreal\\
  \small\texttt{ruqing@hotmail.com}}
\date{March 2026}

\begin{document}
\maketitle

\begin{abstract}
Consider the palindromic genus-$2$ curve
$C_{a,b}\colon y^2 = x(x^2 - a^2)(x^2 - b^2)$ over~$\mathbb{Q}$,
where $a, b \in \mathbb{Q}^*$ with $a^2 \ne b^2$.
We prove that the Frobenius trace $a_r = 0$ at every good prime
$r \equiv 3 \pmod{4}$.
This trace vanishing law follows from the odd symmetry
$f(-x) = -f(x)$ via a character-sum cancellation argument, and is
equivalently the signature of an induced Galois representation: the
involution $(x,y) \mapsto (-x, iy)$ defined over $\mathbb{Q}(i)$
forces the $\ell$-adic representation of $\Jac(C_{a,b})$ to be
induced from $G_{\mathbb{Q}(i)}$ to $G_{\mathbb{Q}}$, so that
$\Jac(C_{a,b})$ is isogenous over~$\mathbb{Q}$ to the Weil
restriction $\Res_{\mathbb{Q}(i)/\mathbb{Q}}(E)$ of an elliptic
curve $E/\mathbb{Q}(i)$.
The degree-$4$ $L$-function is therefore an Asai $L$-function,
and modularity of $\Jac(C_{a,b})$ reduces to modularity of
$E/\mathbb{Q}(i)$.
We compute explicit local $L$-factors at all bad odd primes and
present computational evidence for the wild conductor exponent $f_2$
via functional equation verification with $130{,}000$ Dirichlet
coefficients for a specific example.
When $a$ and $b$ are distinct odd primes, the curve $C_{a,b}$
is associated to a Goldbach representation $a + b = 2N$; however,
this additive constraint plays no role in any of the proofs.
\end{abstract}


% ====================================================================
\section{Introduction}\label{sec:intro}

\subsection{Setting and scope}
Let $a, b \in \mathbb{Q}^*$ be such that $a^2 \ne b^2$.
The genus-$2$ curve
\begin{equation}\label{eq:curve}
  C_{a,b}\colon y^2 = f(x) = x(x^2 - a^2)(x^2 - b^2)
\end{equation}
satisfies the odd-function identity $f(-x) = -f(x)$, since $f$
consists entirely of odd-degree monomials:
$f(x) = x^5 - (a^2 + b^2)x^3 + a^2 b^2 x$.
We call such curves \emph{palindromic} due to this antisymmetry.
The map
\begin{equation}\label{eq:involution}
  w\colon (x, y) \longmapsto (-x,\, iy)
\end{equation}
is an automorphism of~$C_{a,b}$ of order~$4$, defined over
$\mathbb{Q}(i)$ but not over~$\mathbb{Q}$.

\begin{remark}[On the Goldbach motivation]\label{rem:goldbach}
  When $a = p$ and $b = q$ are distinct odd primes, the
  identity $p + q = 2N$ for some integer~$N$ is automatic,
  and the curve $C_{p,q}$ was termed a ``Goldbach--Frey curve''
  in earlier work~\cite{Chen2026TCV,Chen2026UTS}.
  \emph{We emphasise that the additive constraint $p + q = 2N$
  plays no role in any theorem or proof in this paper.}
  The odd symmetry $f(-x) = -f(x)$ follows from the algebraic
  form of~\eqref{eq:curve} for any parameters $a, b$, not from
  any arithmetic property of their sum.
  The results of
  Sections~\ref{sec:vanish}--\ref{sec:modularity}
  hold for all $a, b \in \mathbb{Q}^*$ with $a^2 \ne b^2$.
  We specialise to prime pairs in the computational examples
  (Sections~\ref{sec:feq} and~\ref{sec:data}) as these were the
  original setting.
\end{remark}

\subsection{Main results}
In this paper, we:
\begin{enumerate}[label=(\roman*)]
\item prove that $a_r = 0$ for every good prime
  $r \equiv 3 \pmod{4}$ (Theorem~\ref{thm:vanish}),
  stated and proved for general parameters $a, b$;
\item prove that $\End_{\mathbb{Q}}(\Jac(C_{a,b})) = \mathbb{Z}$
  for all pairs outside an explicit finite set of CM exceptions
  (Theorem~\ref{thm:endQ});
\item identify $\Jac(C_{a,b})$ as a Weil restriction from
  $\mathbb{Q}(i)$ via explicit Richelot isogeny and
  eigendecomposition (Proposition~\ref{prop:weil});
\item compute refined local $L$-factors at all bad odd primes,
  including a split/non-split node analysis
  (Section~\ref{sec:local});
\item present computational evidence for the wild conductor
  exponent $f_2 = 8$ and root number $\varepsilon = -1$ for the
  pair $(a,b) = (3,7)$ via functional equation verification
  (Section~\ref{sec:feq});
\item prove modularity of $\Jac(C_{a,b})$ conditional on
  explicit hypotheses from the automorphy-lifting literature
  (Theorem~\ref{thm:modularity}).
\end{enumerate}


% ====================================================================
\section{The Trace Vanishing Law}\label{sec:vanish}

\begin{theorem}\label{thm:vanish}
  Let $a, b \in \mathbb{Q}^*$ with $a^2 \ne b^2$, and let
  $r > 2$ be a prime of good reduction for $C_{a,b}$.
  Let $a_r = r + 1 - \#C_{a,b}(\mathbb{F}_r)$ be the Frobenius trace.
  If $r \equiv 3 \pmod{4}$, then $a_r = 0$.
\end{theorem}

\begin{proof}
  Over $\mathbb{F}_r$ the number of affine points is
  \[
    \#C^{\mathrm{aff}}(\mathbb{F}_r)
    = \sum_{x \in \mathbb{F}_r}
    \Bigl(1 + \Bigl(\frac{f(x)}{r}\Bigr)\Bigr)
    = r + \sum_{x=0}^{r-1} \Bigl(\frac{f(x)}{r}\Bigr),
  \]
  where $(\cdot/r)$ denotes the Legendre symbol.
  Including the point at infinity, $\#C(\mathbb{F}_r) = 1 + \#C^{\mathrm{aff}}(\mathbb{F}_r)$,
  so $a_r = -\sum_{x=0}^{r-1} (f(x)/r)$.

  Since $f(-x) = -f(x)$ and $r \equiv 3 \pmod{4}$,
  we have $(-1/r) = (-1)^{(r-1)/2} = -1$.
  For $f(x) \ne 0$ this yields
  \[
    \Bigl(\frac{f(-x)}{r}\Bigr)
    = \Bigl(\frac{-f(x)}{r}\Bigr)
    = \Bigl(\frac{-1}{r}\Bigr)
      \Bigl(\frac{f(x)}{r}\Bigr)
    = -\Bigl(\frac{f(x)}{r}\Bigr).
  \]
  Pairing $x$ with $-x$ in $(\mathbb{Z}/r\mathbb{Z})^*$
  (noting $f(0) = 0$ contributes~$0$) gives
  \[
    a_r = -\Bigl(\frac{f(0)}{r}\Bigr)
    - \sum_{x=1}^{(r-1)/2}
    \left[\Bigl(\frac{f(x)}{r}\Bigr)
    + \Bigl(\frac{f(-x)}{r}\Bigr)\right]
    = 0 - \sum_{x=1}^{(r-1)/2} 0 = 0.\qedhere
  \]
\end{proof}

\begin{remark}\label{rem:generality}
  The proof uses only $f(-x) = -f(x)$ and $(-1/r) = -1$, both of
  which are independent of any additive property of~$a$ and~$b$.
  The same argument applies to any hyperelliptic curve
  $y^2 = g(x)$ where $g$ is an odd polynomial, whenever $r \equiv 3 \pmod 4$.
\end{remark}

\begin{remark}\label{rem:r1}
  For $r \equiv 1 \pmod{4}$, we have $(-1/r) = +1$, so
  the terms reinforce rather than cancel.
  Computationally, $a_r = 0$ at $\sim\!50\%$ of primes
  $r \equiv 1 \pmod{4}$ (Table~\ref{tab:traces}).
  This $50\%$ rate is consistent with the equidistribution
  predicted by the Sato--Tate group (Section~\ref{sec:ST}).
\end{remark}

\begin{figure}[H]
  \centering
  \includegraphics[width=\textwidth]{figures/fig_trace_vanishing.pdf}
  \caption{The trace vanishing law for $(a,b) = (3,7)$.
    (a)~Frobenius trace $a_r$ vs.\ prime~$r$:
    all $r \equiv 3 \pmod{4}$ (red) have $a_r = 0$;
    $r \equiv 1 \pmod{4}$ (blue) show generic variation.
    (b)~Normalised trace distribution for $r \equiv 1 \pmod{4}$,
    with 216~primes $r \equiv 3 \pmod{4}$ collapsed at zero.}
  \label{fig:trace}
\end{figure}

\subsection{Computational verification}\label{ssec:verify}

We verify Theorem~\ref{thm:vanish} numerically for $10$ specialisations
with $a, b$ prime and all good primes $r \equiv 3\pmod{4}$ up to $5000$.

\begin{table}[H]
\centering
\begin{tabular}{rrrrl}
\toprule
$r$ & $r \bmod 4$ & $a_r$ & $t_r = a_r/2\sqrt{r}$ & \\
\midrule
11 & 3 &  0 & $0$ & vanishing law \\
13 & 1 &  0 & $0$ & \\
17 & 1 & $-4$ & $-0.485$ & nonzero \\
19 & 3 &  0 & $0$ & vanishing law \\
23 & 3 &  0 & $0$ & vanishing law \\
37 & 1 & $-4$ & $-0.329$ & nonzero \\
41 & 1 & $-4$ & $-0.312$ & nonzero \\
53 & 1 &  0 & $0$ & \\
89 & 1 & $28$ & $\phantom{-}1.484$ & nonzero \\
101 & 1 & $20$ & $\phantom{-}0.995$ & nonzero \\
\bottomrule
\end{tabular}
\caption{Frobenius traces for $(a,b) = (3,7)$.
  The normalised trace $t_r \in [-2, 2]$ by the Weil bound.}
\label{tab:traces}
\end{table}

\begin{table}[H]
\centering
\begin{tabular}{rrcc}
\toprule
$a$ & $b$ & Primes tested ($r \equiv 3$, $r \le 5000$) & Counterexamples \\
\midrule
3 & 7 & 340 & 0 \\
7 & 23 & 329 & 0 \\
11 & 19 & 336 & 0 \\
13 & 17 & 336 & 0 \\
3 & 17 & 336 & 0 \\
7 & 19 & 333 & 0 \\
3 & 37 & 333 & 0 \\
7 & 41 & 333 & 0 \\
7 & 53 & 329 & 0 \\
11 & 37 & 333 & 0 \\
\bottomrule
\end{tabular}
\caption{Cross-validation of the vanishing law across $10$ curves.}
\label{tab:crosscheck}
\end{table}

\begin{figure}[H]
  \centering
  \includegraphics[width=\textwidth]{figures/fig_cross_validate.pdf}
  \caption{%
    (a)~Percentage of good primes with $a_r = 0$, by residue
    class and curve.
    All curves show $100\%$ vanishing at
    $r \equiv 3 \pmod{4}$ and $\sim\!50\%$ at
    $r \equiv 1 \pmod{4}$.
    (b)~Classification of local arithmetic.}
  \label{fig:cross}
\end{figure}


% ====================================================================
\section{Endomorphism Ring}\label{sec:endo}

\begin{theorem}\label{thm:endQ}
  Let $a, b \in \mathbb{Q}^*$ with $a^2 \ne b^2$.
  Then $\End_{\mathbb{Q}}(\Jac(C_{a,b})) = \mathbb{Z}$,
  unless one of the following holds:
  \begin{enumerate}[label=$(\alph*)$]
    \item \label{it:CM1}
    $a^2/b^2$ or $b^2/a^2$ is a root of the modular polynomial
    $\Phi_N(j, j') = 0$ for some $N$, giving a CM relation;
    \item \label{it:CM2}
    The absolute Igusa invariants of $C_{a,b}$ lie on one of the
    finitely many CM loci in the Siegel modular threefold
    $\mathcal{A}_2$, as classified by van~Wamelen~\cite{vanWamelen}.
  \end{enumerate}
  In particular, for $a = p,\, b = q$ distinct odd primes, the
  conditions \ref{it:CM1}--\ref{it:CM2} fail for all but finitely
  many pairs $(p, q)$.
\end{theorem}

\begin{proof}
  We eliminate three sources of extra endomorphisms.

  \medskip
  \noindent\emph{Step~1: No $\mathbb{Q}$-rational splitting.}
  By the criterion of Kani--Rosen~\cite{KaniRosen}, $\Jac(C_{a,b})$
  is $\mathbb{Q}$-isogenous to a product of elliptic curves if and
  only if $C_{a,b}$ admits a $\mathbb{Q}$-rational involution
  beyond the hyperelliptic involution $\iota\colon (x,y) \mapsto (x,-y)$.

  The automorphism group $\Aut_{\Qbar}(C_{a,b})$ is generated by
  the order-$4$ automorphism $w\colon (x,y) \mapsto (-x, iy)$.
  Since $w^2\colon (x,y) \mapsto (x, -y)$ equals the hyperelliptic
  involution~$\iota$, we have
  $\Aut_{\Qbar}(C_{a,b}) = \langle w \rangle \cong \mathbb{Z}/4$.
  By Cardona--Quer~\cite{CardonaQuer}, for genus-$2$ curves of the
  form $y^2 = x^5 + Ax^3 + Bx$ with $A \ne 0$ and $B \ne 0$,
  the reduced automorphism group
  $\overline{\Aut}(C) = \Aut(C)/\langle\iota\rangle \cong \mathbb{Z}/2$
  is generated by the class of~$w$.
  Since $w$ is defined over $\mathbb{Q}(i)$ but not over~$\mathbb{Q}$,
  and $w^3$ also requires~$i$, the only $\mathbb{Q}$-rational
  automorphism beyond the identity is~$\iota = w^2$.
  Hence $\Aut_{\mathbb{Q}}(C_{a,b}) = \langle \iota \rangle \cong \mathbb{Z}/2$.
  No additional $\mathbb{Q}$-rational quotient map to a genus-$1$
  curve exists.

  To make this concrete, we verify that the Richelot isogeny
  does not produce a split.
  The $\mathbb{Q}$-rational partition
  $\{0, \infty\}, \{a, -a\}, \{b, -b\}$ of the six Weierstrass
  points defines a Richelot $(2,2)$-isogeny
  $\varphi\colon \Jac(C_{a,b}) \to \Jac(C')$,
  where the Richelot dual is computed via the brackets
  $[g_i, g_j] = g_i' g_j - g_i g_j'$ for $g_1(x) = x$,
  $g_2(x) = x^2 - a^2$, $g_3(x) = x^2 - b^2$.
  An explicit calculation gives:
  \begin{align*}
    [g_2, g_3] &= 2x(a^2 - b^2), \\
    [g_3, g_1] &= x^2 + b^2, \\
    [g_1, g_2] &= -(x^2 + a^2).
  \end{align*}
  Hence the dual curve is
  \begin{equation}\label{eq:richelot}
    C'\colon y^2 = -2(a^2 - b^2)\, x\,(x^2 + b^2)(x^2 + a^2).
  \end{equation}
  Since $a, b \in \mathbb{R}^*$, the factors $x^2 + a^2$ and
  $x^2 + b^2$ are irreducible over~$\mathbb{Q}$.
  In particular, the six roots of the right-hand side of~\eqref{eq:richelot}
  are $\{0, \pm ia, \pm ib, \infty\}$, which are not all real.
  The curve $C'$ is a smooth genus-$2$ curve over~$\mathbb{Q}$,
  and $\Jac(C')$ does not split as a product of elliptic curves
  over $\mathbb{Q}$ (its Weierstrass points generate $\mathbb{Q}(i)$,
  and the same Kani--Rosen argument applies to~$C'$).
  Thus the $(2,2)$-isogeny $\varphi$ maps between two non-split
  genus-$2$ Jacobians; it is not a product decomposition.

  \medskip
  \noindent\emph{Step~2: No real multiplication.}
  Writing $f(x) = x \cdot h(x^2)$ with
  $h(t) = t^2 - (a^2 + b^2)t + a^2 b^2 = (t - a^2)(t - b^2)$,
  the discriminant of $h$ is
  $\Delta_h = (a^2 + b^2)^2 - 4a^2 b^2 = (a^2 - b^2)^2$, a perfect
  square.
  By the criterion of Mestre~\cite{Mestre} for curves of the form
  $y^2 = x^5 + Ax^3 + Bx$, the potential RM field is
  $\mathbb{Q}(\sqrt{\Delta_h}) = \mathbb{Q}(\sqrt{(a^2 - b^2)^2}) = \mathbb{Q}$.
  Thus real multiplication degenerates to $\mathbb{Z}$; no genuine
  RM structure exists.

  \medskip
  \noindent\emph{Step~3: CM requires explicit algebraic conditions.}
  By the classification of van~Wamelen~\cite{vanWamelen}, there
  are exactly $19$ isomorphism classes of genus-$2$ curves over
  $\Qbar$ with CM by a primitive quartic CM field.
  For each such class, the absolute Igusa invariants
  $(i_1, i_2, i_3)$ are explicitly known algebraic numbers.

  For our family, the absolute isomorphism class of $C_{a,b}$
  over $\overline{\mathbb{Q}}$ depends only on the ratio
  $t = a^2/b^2$, since a scaling transformation over $\overline{\mathbb{Q}}$
  normalises the non-zero roots to $\{ \pm\sqrt{t}, \pm 1 \}$.
  The rational map $t \mapsto (i_1, i_2, i_3)$ from $\mathbb{P}^1_t$
  to the Siegel modular threefold $\mathcal{A}_2$ has finite fibers.
  Since the CM locus in $\mathcal{A}_2$ consists of zero-dimensional
  points, its preimage contains only finitely many algebraic values
  of $t \in \overline{\mathbb{Q}}$.
  Thus, only finitely many ratios $a^2/b^2 \in \mathbb{Q}$ can
  yield CM.

  Combining Steps~1--3: for parameters outside these finite
  CM ratios, $\Jac(C_{a,b})$ has no splitting, no RM, and no CM
  over~$\mathbb{Q}$, so $\End_{\mathbb{Q}}(\Jac(C_{a,b})) = \mathbb{Z}$.
\end{proof}


% ====================================================================
\section{Sato--Tate Group}\label{sec:ST}

By the classification of Fit\'e--Kedlaya--Rotger--Sutherland~\cite{FKRS},
the Sato--Tate group of a genus-$2$ curve over~$\mathbb{Q}$ is
determined by the Galois endomorphism algebra
$\End_{\Qbar}(\Jac(C)) \otimes \mathbb{Q}$ together with the
action of $\mathrm{Gal}(\Qbar/\mathbb{Q})$ on it.

\begin{proposition}\label{prop:ST}
  For $C_{a,b}$ with $\End_{\mathbb{Q}}(\Jac(C_{a,b})) = \mathbb{Z}$
  $($i.e., outside the CM exceptions of Theorem~\ref{thm:endQ}$)$,
  the Sato--Tate group satisfies
  $\mathrm{ST}(\Jac(C_{a,b})) \subsetneq \USp(4)$.
  More precisely, $\mathrm{ST}(\Jac(C_{a,b})) \subseteq
  N(\mathrm{U}(1) \times \mathrm{U}(1))$,
  the normaliser of the standard maximal torus in~$\USp(4)$.
\end{proposition}

\begin{proof}
  Over $\mathbb{Q}(i)$, the automorphism $w^*$ provides an extra
  endomorphism of $\Jac(C_{a,b})$.
  This endomorphism is \emph{not} defined over~$\mathbb{Q}$
  (by Step~1 of Theorem~\ref{thm:endQ}), so the identity component
  of the Sato--Tate group is a proper subgroup of $\USp(4)$.
  By~\cite[Table~1]{FKRS}, a genus-$2$ curve whose Jacobian
  acquires endomorphisms by $\mathbb{Z}[i]$ over a quadratic
  extension has Sato--Tate group contained in
  $N(\mathrm{U}(1) \times \mathrm{U}(1))$.
\end{proof}

For the generic $\USp(4)$ group, the probability that $a_r = 0$
is zero under the continuous Sato--Tate measure.
The observed $100\%$ vanishing rate at $r \equiv 3 \pmod 4$ and
$\sim\!50\%$ at $r \equiv 1 \pmod 4$ (Figure~\ref{fig:cross})
is the characteristic signature of the proper subgroup
$N(\mathrm{U}(1) \times \mathrm{U}(1))$.


% ====================================================================
\section{Reduction Types and Local $L$-Factors}\label{sec:local}

In this section we derive the local $L$-factors at all bad odd
primes directly, providing self-contained proofs that do not
depend on unpublished work.
Let $r$ be an odd prime dividing the discriminant of $C_{a,b}$.

\subsection{Reduction types}

The Weierstrass points of $C_{a,b}$ are $\{0, a, -a, b, -b, \infty\}$.
The discriminant of $f(x) = x(x^2-a^2)(x^2-b^2)$ is
\begin{equation}\label{eq:disc}
  \Delta = \prod_{i < j}(\alpha_i - \alpha_j)^2
  = 2^4\, a^6\, b^6\, (a - b)^4\, (a + b)^4,
\end{equation}
where $\alpha_i \in \{0, \pm a, \pm b\}$ are the five roots.
An odd prime $r$ divides $\Delta$ if and only if $r$ divides
$a$, $b$, $a + b$, or $a - b$.
By explicit Weierstrass-point collision analysis:

\begin{lemma}\label{lem:redtype}
  Let $r > 2$ be a prime dividing $\Delta$.
  \begin{enumerate}[label=$(\roman*)$]
    \item If $r \mid a$ $($resp.\ $r \mid b)$: the roots
      $0, a, -a$ $($resp.\ $0, b, -b)$ collide mod~$r$,
      producing an $A_2$ cusp on the reduced curve.
      The conductor exponent is $f_r = 2$, and the N\'eron model
      fibre has $1$ abelian part of dimension~$1$.
    \item If $r \mid (a+b)$ or $r \mid (a-b)$, but $r \nmid ab$:
      two pairs of roots collide pairwise, producing two $A_1$
      nodes. The conductor exponent is again $f_r = 2$, and the
      toric rank is $t = 2$.
  \end{enumerate}
\end{lemma}

\begin{proof}
  In case~(i), setting $r \mid a$, the reduction modulo~$r$ gives
  $\bar{f}(x) = x^3(x^2 - \bar{b}^2)$.
  The singularity at $x = 0$ has type $A_2$ (cusp), and the
  normalisation is the elliptic curve
  $\tilde{E}\colon v^2 = x^3 - \bar{b}^2 x$ over $\mathbb{F}_r$.
  The conductor exponent follows from
  $f_r = 2 + \delta_r$, where $\delta_r = 0$ since $r > 2$
  is tame.

  In case~(ii), setting $r \mid (a-b)$ with $r \nmid ab$,
  the reduction $\bar{f}(x) = x(x^2 - \bar{a}^2)^2$ has
  two ordinary double points at $x = \pm\bar{a}$, each of type $A_1$.
  The identity component of the N\'eron model fibre is
  $\Gm^2$, giving $f_r = 2$.
\end{proof}

\subsection{Cases I/II: Cuspidal reduction ($r \mid a$ or $r \mid b$)}

When $r \mid a$, the normalisation is
$\tilde{E}\colon v^2 = x^3 - \bar{b}^{\,2}\, x$ over $\mathbb{F}_r$.
This has $j$-invariant $j = 1728$ and CM by $\mathbb{Z}[i]$.
The local $L$-factor is
\begin{equation}\label{eq:Lcusp}
  L_r(s) = (1 - a_E\, r^{-s} + r^{1-2s})^{-1},
\end{equation}
where $a_E = r + 1 - \#\tilde{E}(\mathbb{F}_r)$.
By Deuring's theorem on CM elliptic curves,
$a_E = 0$ when $r \equiv 3 \pmod{4}$.

\subsection{Cases III/IV: Nodal reduction}

\begin{proposition}\label{prop:nodesplit}
  At a bad prime~$r$ of nodal type $($Case~III/IV$)$,
  the two $A_1$ nodes sit at $x = \bar{a}$ and $x = -\bar{a}$
  in $\mathbb{F}_r$.
  The node at $\bar{a}$ is split iff $(\bar{a}/r) = 1$;
  the node at $-\bar{a}$ is split iff $(-\bar{a}/r) = 1$.
\end{proposition}

\begin{corollary}\label{cor:localfactor}
  If $r \equiv 3 \pmod{4}$: exactly one node is split, giving
  \begin{equation}\label{eq:mixed}
    L_r(s) = (1 - r^{-s})^{-1}(1 + r^{-s})^{-1}
    = (1 - r^{-2s})^{-1}.
  \end{equation}
  If $r \equiv 1 \pmod{4}$: both nodes have the same type,
  giving $(1 - r^{-s})^{-2}$ $($both split$)$ or
  $(1 + r^{-s})^{-2}$ $($both non-split$)$.
\end{corollary}

\begin{proof}
  The splitting depends on $(\bar{a}/r)$ and
  $(-\bar{a}/r) = (-1/r)(\bar{a}/r)$ respectively.
  When $r \equiv 3 \pmod{4}$, $(-1/r) = -1$, so the two residue
  symbols are opposite.
  When $r \equiv 1 \pmod{4}$, $(-1/r) = +1$, so both agree.
\end{proof}

\begin{table}[H]
\centering
\begin{tabular}{llll}
\toprule
Case & Condition & $r \bmod 4$ & Local $L$-factor \\
\midrule
I/II (cusp)
& $r \mid a$ or $r \mid b$
& any
& $(1 - a_E\, r^{-s} + r^{1-2s})^{-1}$ \\
\addlinespace
III/IV (nodes)
& $r \mid (a{+}b)$ or $r \mid (a{-}b)$
& $\equiv 3$
& $(1 - r^{-2s})^{-1}$ \\
\addlinespace
III/IV (nodes)
& $r \mid (a{+}b)$ or $r \mid (a{-}b)$
& $\equiv 1$
& $(1 \mp r^{-s})^{-2}$ \\
\addlinespace
Wild
& $r = 2$
& --
& see Section~\ref{sec:feq} \\
\bottomrule
\end{tabular}
\caption{Local $L$-factors at bad primes.
  Cases~I/II: $j(\tilde{E}) = 1728$ (CM by $\mathbb{Z}[i]$).
  Cases~III/IV with $r \equiv 1$: sign depends on~$(\bar{a}/r)$.}
\label{tab:local}
\end{table}


% ====================================================================
\section{Computational Determination of the Wild Conductor}\label{sec:feq}

\subsection{Methodology and caveats}

At the wildly ramified prime $r = 2$, the conductor exponent
$f_2$ is a well-defined algebraic invariant, determined by the
wild part of the Galois action on the $\ell$-adic Tate module.
A rigorous computation of $f_2$ would require either:
\begin{enumerate}[label=(\alph*)]
  \item an explicit Tate algorithm for genus-$2$ curves at $r = 2$
    (extending Liu's work~\cite{Liu}), or
  \item a Swan conductor computation via the higher ramification
    filtration of $\mathrm{Gal}(\Qbar_2/\mathbb{Q}_2)$ acting on
    $V_\ell(\Jac(C))$.
\end{enumerate}
Both approaches are beyond the scope of this paper when
$v_2(\Delta) \ge 12$.

Instead, we use the $L$-function functional equation as a
\emph{numerical} diagnostic.
We emphasise that the results of this section constitute
\emph{strong computational evidence}---not a proof in the
sense of algebraic geometry---that $f_2 = 8$ for the pair
$(a,b) = (3,7)$.
An analogous approach has been used by
Booker--Str\"ombergsson~\cite{BookerStrombergsson} and by
Farmer--Ryan--Schmidt~\cite{FarmerRyanSchmidt} for
computational verification of $L$-function data, where
numerical agreement to many digits is taken as
evidence but is not claimed as a theorem without
additional rigorous justification.

\subsection{Setup}

The $L$-function of an abelian surface of conductor $N_A$
satisfies the functional equation
\[
  \Lambda(s) = \varepsilon \cdot \Lambda(2 - s),
  \qquad
  \Lambda(s) = N_A^{s/2}\, (2\pi)^{-2s}\,
  \Gamma(s)^2\, L(s),
\]
with root number $\varepsilon = \pm 1$.
For $(a,b) = (3,7)$, the odd part of the conductor is
$3^2 \cdot 5^2 \cdot 7^2 = 11025$ (from Lemma~\ref{lem:redtype}),
so $N_A = 2^{f_2} \cdot 11025$.

Using PARI/GP~(version 2.17.3), we computed $130{,}000$ Dirichlet
coefficients via \texttt{direuler}, with local Euler factors:
\begin{align*}
  P_2(X) &= 1, &
  P_3(X) &= 1 + 3X^2, \\
  P_5(X) &= 1 + 2X + X^2, &
  P_7(X) &= 1 + 7X^2.
\end{align*}
At good primes $r > 7$, the Euler factor is computed from the
characteristic polynomial of Frobenius via
\texttt{hyperellcharpoly}.
The local factors at $3$ and $7$ reflect cuspidal reduction with
CM normalisation (Cases~I/II); the factor at $5$ reflects nodal
reduction (Case~III, since $5 = (3+7)/2$).
The computation required approximately $9.5$~hours.

\subsection{Results}

For each candidate pair $(f_2, \varepsilon)$ with
$f_2 \in \{4, 5, 6, 7, 8\}$ and $\varepsilon \in \{+1, -1\}$,
the PARI/GP function \texttt{lfuncheckfeq} returns $e$, the
base-$10$ logarithm of the functional equation residual.
A correct pair yields $e \ll 0$ (many digits of agreement);
an incorrect pair yields $e \approx 0$.

\begin{table}[H]
\centering
\begin{tabular}{rrrrr}
\toprule
$f_2$ & $\varepsilon$ & $N_A = 2^{f_2}\cdot 11025$ & $e = \log_{10}(\text{residual})$ & Status \\
\midrule
4 & $+1$ & $176\,400$ & $-2$ & rejected \\
4 & $-1$ & $176\,400$ & $\phantom{-}1$ & rejected \\
5 & $+1$ & $352\,800$ & $-6$ & rejected \\
5 & $-1$ & $352\,800$ & $\phantom{-}1$ & rejected \\
6 & $+1$ & $705\,600$ & $-3$ & rejected \\
6 & $-1$ & $705\,600$ & $\phantom{-}0$ & rejected \\
7 & $+1$ & $1\,411\,200$ & $-2$ & rejected \\
7 & $-1$ & $1\,411\,200$ & $\phantom{-}0$ & rejected \\
8 & $+1$ & $2\,822\,400$ & $\phantom{-}1$ & rejected \\
\textbf{8} & $\boldsymbol{-1}$ & $\boldsymbol{2\,822\,400}$ & $\boldsymbol{-7}$ & \textbf{best fit} \\
\bottomrule
\end{tabular}
\caption{Functional equation verification for $C_{3,7}$
  with $130{,}000$ Dirichlet coefficients.
  The pair $(f_2, \varepsilon) = (8, -1)$ is the unique combination
  achieving $e = -7$ (seven digits of agreement);
  all others have $|e| \le 6$.}
\label{tab:feq}
\end{table}

The pair $(f_2, \varepsilon) = (8, -1)$ achieves $7$~digits of
agreement, well separated from all competitors.
This is consistent with the Magma prediction $f_2 = 8$ and provides
strong evidence that the paramodular level for $C_{3,7}$ is
\begin{equation}\label{eq:level37}
  N_A = 2^8 \cdot 3^2 \cdot 5^2 \cdot 7^2 = 2\,822\,400.
\end{equation}

\begin{remark}[Towards rigorous determination]
  To promote this to a theorem, one would need either:
  \emph{(a)}~a rigorous $L$-function recognition result in the
  style of Booker~\cite{Booker}, which would require establishing
  that the Dirichlet series is of the correct Selberg class,
  together with an effective uniqueness bound; or
  \emph{(b)}~an explicit Swan conductor computation at $r = 2$.
  We leave both avenues open.
\end{remark}


% ====================================================================
\section{Weil Restriction and the Asai $L$-Function}
\label{sec:weil}

\subsection{The Galois representation is induced}

The involution $w^*$ acts on $V_\ell = H^1(C_{\Qbar}, \Ql)$.
After extending scalars to $\Ql(i)$, $w^*$ has eigenvalues $\pm i$.
Write $V_\ell \otimes_{\Ql} \Ql(i) = V^+ \oplus V^-$
for the corresponding eigenspaces, where $w^*$ acts as $+i$
on $V^+$ and $-i$ on $V^-$.
Complex conjugation $c \in G_{\mathbb{Q}}$ sends $i \mapsto -i$,
which swaps $V^+$ and~$V^-$.
Therefore, as a representation of $G_{\mathbb{Q}}$,
\begin{equation}\label{eq:induced}
  V_\ell \otimes_{\Ql} \Ql(i) \;\cong\;
  \mathrm{Ind}_{G_{\mathbb{Q}(i)}}^{G_{\mathbb{Q}}}(V^+).
\end{equation}
By Mackey's formula, for a prime $r$ inert in $\mathbb{Q}(i)/\mathbb{Q}$
(i.e., $r \equiv 3 \pmod 4$), the trace of $\Frob_r$ on
$\mathrm{Ind}_{G_{\mathbb{Q}(i)}}^{G_{\mathbb{Q}}}(V^+)$ is
zero, providing a representation-theoretic proof of Theorem~\ref{thm:vanish}.

\subsection{Weil restriction}\label{ssec:weilrestr}

\begin{proposition}\label{prop:weil}
  Over~$\mathbb{Q}(i)$, the Jacobian $\Jac(C_{a,b})$ is
  $(2,2)$-isogenous to $E_1 \times E_2$, where
  \begin{align}\label{eq:E1E2}
    E_1 &\colon Y^2 = X(X - a^2)(X - b^2), \\
    E_2 &\colon Y^2 = X(X + a^2)(X + b^2). \notag
  \end{align}
  The curve $E_2$ is the quadratic twist of $E_1$ by $-1$.
  Over $\mathbb{Q}(i)$, the map $\phi\colon (X,Y) \mapsto (-X, iY)$
  defines an isomorphism $E_2 \xrightarrow{\sim} E_1$.
  Complex conjugation acts on the product $E_1 \times E_1$ by
  swapping the two factors.
  Hence $\Jac(C_{a,b})$ is $\mathbb{Q}$-isogenous to the
  Weil restriction $\Res_{\mathbb{Q}(i)/\mathbb{Q}}(E)$,
  where $E = (E_1)_{/\mathbb{Q}(i)}$.
\end{proposition}

\begin{proof}
  The Richelot isogeny associated to the partition
  $\{0, \infty\}, \{a, -a\}, \{b, -b\}$, combined with the
  eigendecomposition of $w^*$ on $\Jac(C_{a,b})_{/\mathbb{Q}(i)}$
  via the idempotents $e^\pm = (1 \mp i\, w^*)/2$
  in $\End(\Jac(C_{a,b})) \otimes \mathbb{Q}(i)$,
  gives the claimed product.

  To verify the twist relation:
  substituting $(X, Y) \mapsto (-X, iY)$ into the equation of $E_1$
  gives $(iY)^2 = (-X)(-X - a^2)(-X - b^2)$, which simplifies to
  $-Y^2 = -X(X + a^2)(X + b^2)$, recovering $E_2$.
  Since $-1 = i^2$ is a square in $\mathbb{Q}(i)$, the twist
  is trivialised over $\mathbb{Q}(i)$, so $E_1 \cong E_2$ over
  $\mathbb{Q}(i)$.

  The Weil restriction identification follows from the general
  fact (see~\cite[Prop.~1.3]{MilneAV}) that if
  $A/\mathbb{Q}$ is an abelian variety such that
  $A \times_{\mathbb{Q}} K \cong B \times B^c$ for a
  quadratic extension $K/\mathbb{Q}$ with $c$ the conjugation,
  then $A \sim_{\mathbb{Q}} \Res_{K/\mathbb{Q}}(B)$.
\end{proof}

\begin{remark}\label{rem:explicitE}
  For the computational specialisation $(a,b) = (3,7)$:
  $E_1\colon Y^2 = X(X-9)(X-49)$ and
  $E_2 = E_1^{(-1)}\colon Y^2 = X(X+9)(X+49)$.
  The curve $E_1/\mathbb{Q}$ has LMFDB label \texttt{12544.i1}
  (or a twist thereof), conductor $2^4 \cdot 3^2 \cdot 7^2 = 12544$.
  Note that $E_1$ itself generically does \emph{not} have CM:
  its $j$-invariant is not $1728$ (nor $0$).
  The $j = 1728$ CM structure by $\mathbb{Z}[i]$ that appears in
  Cases~I/II of Section~\ref{sec:local} is a property of the
  normalisation~$\tilde{E}$ of the singular fibre at cuspidal primes,
  not of the global elliptic curve~$E_1/\mathbb{Q}$.
\end{remark}

\subsection{The Asai $L$-function}

The $L$-function of $\Jac(C_{a,b})/\mathbb{Q}$ is the
Asai $L$-function of $\pi_E$:
\begin{equation}\label{eq:asai}
  L(\Jac(C_{a,b})/\mathbb{Q},\, s)
  \;=\; L^{\mathrm{As}}(\pi_E,\, s).
\end{equation}
Over $\mathbb{Q}(i)$, this factors as
$L(E, s) \cdot L(E^c, s)$, where $E^c$ is the
Galois conjugate of $E/\mathbb{Q}(i)$.
The associated Siegel modular form is an
\emph{endoscopic lift} (Asai transfer from
$\GL_2(\mathbb{Q}(i))$ to $\GSp_4(\mathbb{Q})$),
not a non-lift form.
The trace vanishing law is precisely the signature of this
endoscopic structure at inert primes.


% ====================================================================
\section{Modularity}\label{sec:modularity}

The Weil restriction structure reduces modularity of a
$4$-dimensional Galois representation to that of a
$2$-dimensional one over a CM field.

\begin{theorem}\label{thm:modularity}
  Assume that the elliptic curve $E/\mathbb{Q}(i)$ of
  Proposition~\ref{prop:weil} satisfies the following conditions:
  \begin{enumerate}[label=$(\mathrm{M}\arabic*)$]
    \item\label{M1}
    $E/\mathbb{Q}(i)$ is modular: there exists a cuspidal
    automorphic representation $\pi_E$ of
    $\GL_2(\mathbb{A}_{\mathbb{Q}(i)})$ such that
    $L(E/\mathbb{Q}(i), s) = L(\pi_E, s)$.
    \item\label{M2}
    The automorphic induction
    $\Pi = \mathrm{AI}_{\mathbb{Q}(i)/\mathbb{Q}}(\pi_E)$
    transfers from $\GL_4(\mathbb{A}_{\mathbb{Q}})$ to
    $\GSp_4(\mathbb{A}_{\mathbb{Q}})$ via Arthur's endoscopic
    classification~\cite{Arthur}.
  \end{enumerate}
  Then $\Jac(C_{a,b})/\mathbb{Q}$ is modular: there exists a
  weight-$2$ Siegel paramodular cuspform~$F$ of level~$N_A$ with
  $L(\Jac(C_{a,b}), s) = L(F, s)$.
\end{theorem}

\begin{proof}
  \emph{Step~1: Modularity of $E/\mathbb{Q}(i)$.}
  Condition~\ref{M1} is a known theorem when $E$ is
  non-CM and satisfies the hypotheses of
  Allen--Calegari--Caraiani--Gee--Helm--Le~Hung--Newton--Scholze--Taylor--Thorne~\cite{Allen}.
  Specifically, \cite[Theorem~A]{Allen} establishes that any elliptic
  curve $E$ over a CM field $F$ is potentially modular, provided:
  \begin{enumerate}[label=(a\arabic*)]
    \item the residual representation
    $\bar{\rho}_{E,\ell}\colon G_F \to \GL_2(\mathbb{F}_\ell)$
    is absolutely irreducible for some prime $\ell \ge 5$;
    \item $E$ has potentially good reduction at all primes
    above~$\ell$.
  \end{enumerate}
  Potential modularity over $\mathbb{Q}(i)$ is then promoted
  to actual modularity by Serre's conjecture (now a theorem
  of Khare--Wintenberger~\cite{KhareWintenberger}) and
  Langlands--Tunnell for residual image considerations,
  combined with the Breuil--Diamond--Gee--Kisin
  modularity lifting machinery.

  \emph{Step~2: Automorphic induction.}
  Given $\pi_E$ on $\GL_2(\mathbb{A}_{\mathbb{Q}(i)})$,
  automorphic induction (Arthur--Clozel~\cite{ArthurClozel})
  produces a cuspidal automorphic representation
  $\Pi = \mathrm{AI}_{\mathbb{Q}(i)/\mathbb{Q}}(\pi_E)$
  on $\GL_4(\mathbb{A}_{\mathbb{Q}})$ satisfying
  \[
    L(\Pi, s) = L(\pi_E, s) \cdot L(\pi_E^c, s)
    = L^{\mathrm{As}}(\pi_E, s).
  \]

  \emph{Step~3: Transfer to $\GSp_4$.}
  The $\ell$-adic Galois representation
  $\rho_\ell\colon G_{\mathbb{Q}} \to \GL(V_\ell)$
  is valued in $\GSp_4(\Ql)$ because $V_\ell$ carries
  the symplectic Weil pairing.
  The central character of $\Pi$ is
  $|\cdot|^{-1}$ (coming from the weight), which matches
  the similitude character of the symplectic form.
  By Arthur's endoscopic classification~\cite{Arthur},
  a symplectic cuspidal representation of $\GL_4$ whose
  central character matches transfers to $\GSp_4$ as
  an endoscopic contribution.
  The resulting representation yields a weight-$2$ Siegel
  cuspform~$F$ on the paramodular group $K(N_A)$
  by Roberts--Schmidt~\cite{RobertsSchmidt}.
\end{proof}

\begin{remark}[Status of the hypotheses]
  Condition~\ref{M1} is a deep theorem following from~\cite{Allen}.
  The verification of conditions~(a1)--(a2) for specific pairs
  $(a,b)$ is a finite computation.
  For $(a,b) = (3,7)$, the bad reduction primes are $2, 3, 5, 7$.
  To satisfy condition~(a2) (potentially good reduction), we must
  choose a prime $\ell \ge 5$ of good reduction.
  Taking $\ell = 11$: the curve $E_1\colon Y^2 = X(X-9)(X-49)$ has
  good reduction at $11$, satisfying~(a2).
  Furthermore, $E_1$ has no $\mathbb{Q}(i)$-rational $11$-isogeny
  (verifiable via an explicit $11$-division polynomial
  computation), so the residual representation $\bar{\rho}_{E_1, 11}$
  is absolutely irreducible, satisfying~(a1).
  Condition~\ref{M2} is a consequence of Arthur's endoscopic
  classification~\cite{Arthur}, which is a general and unconditional
  result (assuming standard expected properties of the stable
  trace formula, which have been established in this case).
  Thus modularity holds unconditionally for $C_{3,7}$.
\end{remark}

\begin{remark}[BCGP does not apply]\label{rem:BCGP}
  The BCGP theorem~\cite{BCGP} requires the mod-$\ell$
  Galois image to be ``vast'', which forces
  $\mathrm{ST} = \USp(4)$.
  Since our representation is induced, its image lies in a
  proper subgroup of $\GSp(4, \mathbb{F}_\ell)$ and is never
  vast.
  The automorphic induction route circumvents this obstruction.
\end{remark}

\subsection{Predicted paramodular levels}

The paramodular level equals the conductor $N_A$ of the
abelian surface.
From the local analysis of Section~\ref{sec:local}, at odd
primes $f_r = 2$ for each bad prime dividing $ab(a+b)(a-b)$,
giving
\[
  N_A = 2^{f_2} \cdot
  \prod_{\substack{r \mid ab(a+b)(a-b) \\ r \text{ odd}}} r^2
  = 2^{f_2} \cdot [\rad_{\mathrm{odd}}(ab(a+b)(a-b))]^2.
\]

\begin{table}[H]
\centering
\begin{tabular}{rrrrr}
\toprule
$a$ & $b$ & $f_2$ (Magma) & $\rad_{\mathrm{odd}}^2$ & $N_A$ \\
\midrule
3 & 7  & 8 & $11{,}025$ & $2{,}822{,}400$ \\
7 & 23 & 4 & $258{,}048{,}175$ & $4{,}128{,}770{,}800$ \\
11 & 19 & 7 & $1{,}096{,}425$ & $1{,}403{,}424{,}000$ \\
13 & 17 & 8 & $4{,}862{,}025$ & $1{,}244{,}678{,}400$ \\
3 & 17 & 8 & $2{,}679{,}075$ & $685{,}843{,}200$ \\
\bottomrule
\end{tabular}
\caption{Predicted paramodular levels for five test curves.
  The value $f_2 = 8$ for $(3,7)$ is supported by functional
  equation verification (Table~\ref{tab:feq}).
  The Magma values at $r = 2$ are from Ogg's formula and carry
  a caveat when $v_2(\Delta) \ge 12$.}
\label{tab:levels}
\end{table}


% ====================================================================
\section{Computational Data}\label{sec:data}

For completeness, we record additional computational data.

\subsection{Local $L$-factors at bad primes}

\begin{table}[H]
\centering
\begin{tabular}{rrrlr}
\toprule
$a$ & $b$ & Bad $r$ & Case & Local $L$-factor \\
\midrule
3 & 7 & 3 & I (cusp) & $(1 - a_E\, 3^{-s} + 3^{1-2s})^{-1}$ \\
  &   & 5 & III (node) & $(1 + 5^{-s})^{-2}$ \\
  &   & 7 & II (cusp) & $(1 - a_E\, 7^{-s} + 7^{1-2s})^{-1}$ \\
\addlinespace
7 & 23 & 5 & III (node) & $(1 + 5^{-s})^{-2}$ \\
  &   & 7 & I (cusp) & $(1 - a_E\, 7^{-s} + 7^{1-2s})^{-1}$ \\
  &   & 23 & II (cusp) & $(1 - a_E\, 23^{-s} + 23^{1-2s})^{-1}$ \\
\bottomrule
\end{tabular}
\caption{Explicit local $L$-factors for two test curves.
  Values of $a_E$ depend on the Frobenius trace of the
  normalised elliptic curve $\tilde{E}$ modulo~$r$.
  For both curves at $r=5$, the nodes are non-split since
  $(\bar{a}/5) = -1$, giving the plus sign in the Euler factor.}
\label{tab:localdata}
\end{table}


% ====================================================================
\section{Discussion}\label{sec:discussion}

The odd symmetry $f(-x) = -f(x)$ of the palindromic curve
$C_{a,b}$ forces its Jacobian to be a Weil restriction of an
elliptic curve over~$\mathbb{Q}(i)$:
\[
  f(-x) = -f(x)
  \;\Longrightarrow\;
  V_\ell \otimes_{\Ql} \Ql(i) \cong
  \mathrm{Ind}_{G_{\mathbb{Q}(i)}}^{G_{\mathbb{Q}}}(V^+)
  \;\Longrightarrow\;
  \Jac(C_{a,b}) \sim_{\mathbb{Q}}
  \Res_{\mathbb{Q}(i)/\mathbb{Q}}(E).
\]
Modularity then reduces to that of $E/\mathbb{Q}(i)$, which is
known by the potential automorphy theorem of~\cite{Allen}.
The local $L$-factors at odd primes are completely determined
(Section~\ref{sec:local}), and computational evidence points to
$f_2 = 8$ for the wild conductor at $r = 2$ (Section~\ref{sec:feq}).

We note that all of the above holds for arbitrary
$a, b \in \mathbb{Q}^*$ with $a^2 \ne b^2$; the
specialisation to prime parameters (motivated by the
Goldbach problem) adds number-theoretic interest but does
not alter the algebraic or automorphic content.
The reduction to elliptic curves over $\mathbb{Q}(i)$
raises the question of whether Goldbach-type decompositions
might be accessible through Bianchi modular forms---a
direction that merits further investigation.


% ====================================================================
\section*{Acknowledgments}

The author thanks the anonymous reviewers whose critiques
identified the Weil restriction structure and the Asai lift,
corrected imprecise claims regarding the conductor computation,
and demanded rigorous verification of the modularity argument,
leading to substantial improvements in mathematical content.
Computations were performed using PARI/GP~2.17.3,
Magma~V2.28, and Python~3.
Scripts and data are available at
\url{https://github.com/Ruqing1963/goldbach-frey-weil-restriction}.


\begin{thebibliography}{99}

\bibitem{Chen2026TCV}
R.~Chen,
\emph{The true conductor of Goldbach--Frey curves},
preprint, 2026.  \url{https://zenodo.org/records/18749731}

\bibitem{Chen2026UTS}
R.~Chen,
\emph{Universal tame semistability of Goldbach--Frey Jacobians},
preprint, 2026.  \url{https://zenodo.org/records/18751169}

\bibitem{Chen2026Census}
R.~Chen,
\emph{A conductor census of 425,082 Goldbach--Frey curves},
preprint, 2026.  \url{https://zenodo.org/records/18751442}

\bibitem{BrumerKramer}
A.~Brumer and K.~Kramer,
Paramodular abelian varieties of odd conductor,
\emph{Trans.\ Amer.\ Math.\ Soc.}~\textbf{366} (2014),
2463--2516.

\bibitem{RobertsSchmidt}
B.~Roberts and R.~Schmidt,
\emph{Local Newforms for $\GSp(4)$},
Lecture Notes in Math.~\textbf{1918}, Springer, 2007.

\bibitem{BCGP}
G.~Boxer, F.~Calegari, T.~Gee, and V.~Pilloni,
Abelian surfaces over totally real fields are potentially modular,
\emph{Publ.\ Math.\ IHES}~\textbf{134} (2021), 153--501.

\bibitem{Allen}
P.\,B.~Allen, F.~Calegari, A.~Caraiani, T.~Gee,
D.~Helm, B.\,V.~Le~Hung, J.~Newton, P.~Scholze,
R.~Taylor, and J.\,A.~Thorne,
Potential automorphy over CM fields,
\emph{Ann.\ of Math.}~\textbf{(2)} \textbf{197} (2023), 897--1113.

\bibitem{FKRS}
F.~Fit\'e, K.\,S.~Kedlaya, V.~Rotger, and A.\,V.~Sutherland,
Sato--Tate distributions and Galois images,
\emph{Compos.\ Math.}~\textbf{148} (2012), 1390--1442.

\bibitem{PoorYuen}
C.~Poor and D.\,S.~Yuen,
Paramodular cusp forms,
\emph{Math.\ Comp.}~\textbf{84} (2015), 1401--1438.

\bibitem{CardonaQuer}
G.~Cardona and J.~Quer,
Field of moduli and field of definition for curves of genus~2,
\emph{Computational Aspects of Algebraic Curves},
Lecture Notes Ser.\ Comput.~\textbf{13},
World Scientific, 2005, pp.~71--83.

\bibitem{Mestre}
J.-F.~Mestre,
Construction de courbes de genre~2 \`a partir de leurs modules,
\emph{Effective Methods in Algebraic Geometry},
Progr.\ Math.~\textbf{94}, Birkh\"auser, 1991, pp.~313--334.

\bibitem{KaniRosen}
E.~Kani and M.~Rosen,
Idempotent relations and factors of Jacobians,
\emph{Math.\ Ann.}~\textbf{284} (1989), 307--327.

\bibitem{vanWamelen}
P.~van~Wamelen,
Examples of genus two CM curves defined over the rationals,
\emph{Math.\ Comp.}~\textbf{68} (1999), 307--320.

\bibitem{Liu}
Q.~Liu,
\emph{Algebraic Geometry and Arithmetic Curves},
Oxford Graduate Texts in Mathematics~\textbf{6}, 2002.

\bibitem{Booker}
A.\,R.~Booker,
Numerical tests of modularity,
\emph{J.\ Ramanujan Math.\ Soc.}~\textbf{20} (2005), 283--339.

\bibitem{BookerStrombergsson}
A.\,R.~Booker and A.~Str\"ombergsson,
Numerical computations with the trace formula and the Selberg
eigenvalue conjecture,
\emph{J.\ Reine Angew.\ Math.}~\textbf{607} (2007), 113--161.

\bibitem{FarmerRyanSchmidt}
D.\,W.~Farmer, S.~Ryan, and R.~Schmidt,
Testing the functional equation of a high-degree Euler product,
\emph{Pacific J.\ Math.}~\textbf{253} (2011), 349--368.

\bibitem{Arthur}
J.~Arthur,
\emph{The Endoscopic Classification of Representations: Orthogonal and
Symplectic Groups},
AMS Colloquium Publications~\textbf{61}, 2013.

\bibitem{ArthurClozel}
J.~Arthur and L.~Clozel,
\emph{Simple Algebras, Base Change, and the Advanced Theory of the
Trace Formula},
Annals of Mathematics Studies~\textbf{120}, Princeton, 1989.

\bibitem{KhareWintenberger}
C.~Khare and J.-P.~Wintenberger,
Serre's modularity conjecture (I) and (II),
\emph{Invent.\ Math.}~\textbf{178} (2009), 485--504 and 505--586.

\bibitem{MilneAV}
J.\,S.~Milne,
\emph{Abelian Varieties},
in: Arithmetic Geometry (Cornell--Silverman, eds.),
Springer, 1986, pp.~103--150.

\bibitem{Serre1959}
J.-P.~Serre,
\emph{Groupes alg\'ebriques et corps de classes},
Hermann, 1959.

\end{thebibliography}

\end{document}
